\documentclass{situationist-preprint}

\title{Selection and Applicability}
\msid{SCI-17-b}
\msfamily{SCI-17}
\msdiscipline{Philosophy of Science}
\mstype{Companion manuscript}
\mskeywords{applicability of mathematics; Wigner; selection; Steiner; indispensability; structure}

\begin{document}
\maketitle

\begin{abstract}
SCI-17-a explains the effectiveness of mathematics in science by selection: scientists choose, from a large stock of structures, those that fit the phenomena. A natural objection is that selection explains too little, since it presupposes that some structures fit and says nothing about why. We distinguish two questions: why the mathematics in use fits, and why any mathematics fits at all. Selection answers the first completely. The second is the question of applicability, which is answered by the instantiation of mathematical structure in physical systems, not by any coincidence. Wigner's puzzle concerns the first question and so is dissolved; the second is a question about the world, not a source of surprise.
\end{abstract}

\section{Two questions}

\citet{Wigner1960} found it unreasonable that mathematical concepts developed for internal reasons describe nature with great accuracy. SCI-17-a argues, by its Theorem~2.5, that the effectiveness is not unreasonable, given Assumption~2.3 (mathematics studies very many structures, most never applied) and Assumption~2.4 (scientists retain the structures that fit and do not count the phenomena that none fits). The point that we find the mathematics we look for was made by \citet{Hamming1980}, who doubted that it was the whole answer.

The doubt can be made precise by distinguishing:
\begin{enumerate}
\item[(Q1)] Why do the structures in scientific use fit the phenomena to which they are applied?
\item[(Q2)] Why does any mathematical structure fit any phenomenon?
\end{enumerate}

\begin{proposition}\label{p:q1}
Given Assumptions~2.3 and~2.4 of SCI-17-a and the fact that some structures fit, selection fully answers (Q1).
\end{proposition}

\begin{proof}
By Assumption~2.4, the structures in use are those retained because they fit, and the domains counted as mathematized are those for which some structure was found to fit. That the structures in use fit is then a consequence of the procedure that put them into use. Nothing about the structures needs to be explained beyond their having been selected.
\end{proof}

Selection presupposes an answer to (Q2) and does not provide one. If no structure fitted any phenomenon, there would be nothing to select. The answer to (Q2) is not a further selection effect. It is that physical systems instantiate structures: a system of masses under gravitation has the structure that the differential equations describe, to a good approximation, and a mathematical description fits when the structure it studies is the structure the system has \citep{Pincock2012}. Definition~2.1 of SCI-17-a builds this into the notion of effectiveness.

\section{Why the residue is not a puzzle}

Wigner's sense of a miracle concerned the proportion of successes, which is the subject of Definition~2.2 of SCI-17-a and the domain of (Q1). Once the proportion is explained, what remains, that the world has structure that mathematics can describe, is not surprising relative to any expectation we had reason to hold. Mathematics is the study of structure in general, so if the world has any structure at all, some branch of mathematics describes it. That the world has stable structure is a condition of its being investigable, and it is not something a theory of the application of mathematics is required to explain.

\begin{objection}
\citet{Steiner1998} argues that physicists have discovered laws by following purely formal analogies, such as extending a formalism because it is mathematically natural, and that the success of this practice suggests that the universe is congenial to human mathematical taste, which selection cannot explain.
\end{objection}

\begin{reply}
Formal analogies are themselves selected. The analogical extensions that failed are not remembered, and those that succeeded were extensions of structures already known to fit, so their success inherits that fit. Selection applies to strategies of discovery as it applies to structures. Where a formal extension succeeds, it does so because the extended structure is instantiated, which is the answer to (Q2), not evidence of a congenial universe.
\end{reply}

Selection is compatible with the indispensability of mathematics to physics, which MTH-16 defends: that the structures in use were selected for fit does not imply that the physics could be stated without them. It is also compatible with an inferential account of application, on which mathematics is applied by mappings between empirical and mathematical structures \citep{BuenoColyvan2011}.

\section{Conclusion}

Selection explains why the mathematics in use fits; instantiation explains why there is mathematics to select. Neither leaves an unreasonable effectiveness to be explained.

\begin{thebibliography}{99}\small
\bibitem[Bueno and Colyvan(2011)]{BuenoColyvan2011} Bueno, O. and Colyvan, M. (2011). An inferential conception of the application of mathematics. \emph{No\^{u}s}, 45(2), 345--374.
\bibitem[Hamming(1980)]{Hamming1980} Hamming, R.~W. (1980). The unreasonable effectiveness of mathematics. \emph{The American Mathematical Monthly}, 87(2), 81--90.
\bibitem[Pincock(2012)]{Pincock2012} Pincock, C. (2012). \emph{Mathematics and Scientific Representation}. Oxford University Press.
\bibitem[Steiner(1998)]{Steiner1998} Steiner, M. (1998). \emph{The Applicability of Mathematics as a Philosophical Problem}. Harvard University Press.
\bibitem[Wigner(1960)]{Wigner1960} Wigner, E.~P. (1960). The unreasonable effectiveness of mathematics in the natural sciences. \emph{Communications on Pure and Applied Mathematics}, 13(1), 1--14.
\end{thebibliography}

\end{document}
